\section{项目背景与赛题理解}

\subsection{项目定位}

\texttt{YH\_rv\_cpu} 的目标不是短期堆叠功能点，而是收敛成一套
``可复现、可验证、可交接、可继续优化'' 的比赛级 CPU 工程基线。
这一定位直接决定了当前工作的三个原则：

\begin{itemize}
    \item 任何对外引用的结果必须有 fresh 验证证据；
    \item 任何优化必须先通过回归门禁，再决定是否保留；
    \item 任何文档更新必须与真实 RTL、脚本和证据路径保持一致。
\end{itemize}

\subsection{赛题答疑对当前方案的影响}

根据 2026-04-09 整理的赛方答疑，当前项目的提交边界可以明确为：

\begin{enumerate}
    \item 初赛以设计说明书、技术文档、交付方案为主，不强制要求完整实板闭环。
    \item 仿真记录可以作为性能与验证证据，因此当前 CoreMark、\texttt{riscv-tests} 和
    Vivado pre-board 结果都具备材料价值。
    \item 不要求机械遵守固定五级流水，也不强制必须完整实现特权架构；因此当前
    \texttt{RV32I + Zicsr} 主线是合理且自洽的。
    \item \texttt{5k LUT} 更偏牵引指标而非唯一硬门槛，设计应重点强调工程完整性、
    架构合理性和验证闭环，而不是为了单一数字破坏主线稳定性。
\end{enumerate}

\subsection{当前交付边界}

为了避免口径漂移，本项目将交付内容分为三层：

\begin{table}[H]
    \centering
    \caption{当前交付边界划分}
    \begin{tabular}{p{0.20\textwidth} p{0.25\textwidth} p{0.43\textwidth}}
        \toprule
        层级 & 当前状态 & 说明 \\
        \midrule
        初赛硬交付 & 已具备 & 设计说明书、技术文档、项目说明、答辩口径 \\
        增强支撑材料 & 已具备 & CoreMark、baseline/full-ui 回归、Vivado impl50、FPGA-like probe \\
        外部阻塞项 & 未完成 & 实体板卡 UART/LED 闭环、最终板级 I/O delay 约束 \\
        \bottomrule
    \end{tabular}
\end{table}

\subsection{冻结基线与活跃验证的区分}

当前文档体系明确区分两类结果：

\begin{itemize}
    \item \textbf{冻结比赛基线}：当前对外统一引用的正式口径，包括 CoreMark short/strict、
    baseline 回归、\texttt{impl50} 与 FPGA-like probe。
    \item \textbf{活跃扩展验证}：用于补全设计可信度的最新验证，例如
    \texttt{rv32/rv64 full-ui} 的 fresh 通过结果，以及 redirect 成本 profile。
\end{itemize}

这种区分可以确保说明书既能展示项目真实进展，又不会把实验态结果误报为正式冻结结论。

\subsection{本说明书的组织方式}

本说明书按照 ``赛题理解 $\rightarrow$ 架构设计 $\rightarrow$ 验证证据 $\rightarrow$
FPGA 原型 $\rightarrow$ 风险与后续计划'' 的顺序展开，强调以下两点：

\begin{itemize}
    \item 说明书的每一项关键结果都可以在仓库中找到对应证据路径；
    \item 文档结构为 LaTeX 模块化布局，后续可以直接补入架构框图、波形截图和 Vivado 报告截图。
\end{itemize}

\clearpage
